\section{Introduction}
\label{sec:intro}

Hierarchical Multi-Label Classification (HMC) is a structured prediction
task where each instance is assigned to one or more nodes in a
predefined class hierarchy---a tree or directed acyclic graph (DAG).
HMC arises naturally in scientific domains: categorizing academic papers
into topic taxonomies (ArXiv, WOS), annotating proteins with Gene
Ontology terms, classifying patents into IPC/CPC hierarchies, and
assigning medical codes (ICD-10) to clinical notes.

A key architectural innovation was introduced by Giunchiglia \&
Lukasiewicz~\cite{giunchiglia2018hmc}: the \textbf{R-matrix} (ancestor
closure matrix), a binary $C \times C$ matrix where $R_{ij}=1$ iff class
$i$ is an ancestor of class $j$. Applied during training as a
consistency loss and at inference as bottom-up max-propagation, the
R-matrix enforces the hierarchical constraint $\forall (i,j):
\text{ancestor}(i,j) \implies p_i \geq p_j$, guaranteeing zero
post-inference hierarchy violations.

Recent work~\cite{sette2026hmctorch} demonstrated that a simple MLP with
R-matrix constraint achieves state-of-the-art results on ArXiv (+13.5
Micro-F1 points over HiAGM~\cite{zhou2020hiagm}), WOS, and AAPD
benchmarks, using frozen transformer embeddings without fine-tuning or
graph neural networks. However, the literature leaves fundamental
questions unanswered:

\begin{enumerate}[leftmargin=*,nosep]
  \item \textbf{Where does the benefit come from?} The R-matrix operates
    at two points---during training (consistency loss) and at inference
    (bottom-up reconciliation). Are both necessary? Which dominates?
  \item \textbf{When does it help?} Under what conditions---hierarchy
    depth, branching factor, training set size, number of classes---does
    the R-matrix provide the largest gains?
  \item \textbf{How does it compare to modern alternatives?} In 2026,
    large language models (GPT-4o, Claude) are routinely used for
    zero-shot classification. Does a simple MLP with hierarchical
    constraints still have a role?
\end{enumerate}

This paper provides the first systematic empirical investigation of
these questions. Our contributions are:

\begin{itemize}[leftmargin=*,nosep]
  \item A \textbf{4-way ablation} isolating training-time consistency
    loss from inference-time reconciliation across three hierarchy
    types (shallow tree, deep tree, DAG). We find that inference-time
    reconciliation alone matches or exceeds the full R-matrix, while
    training-time consistency loss \emph{degrades} performance.
  \item \textbf{Scaling laws for HMC}: the first systematic study of how
    HMC performance scales with number of classes $C$ (50--4,500),
    hierarchy depth $D$ (2--13), and training samples per class
    $N/C$. The R-matrix benefit follows $\Delta\text{F1} \propto
    D \cdot (N/C)^{-0.5}$, providing predictive guidance for when
    hierarchical constraints are worthwhile.
  \item \textbf{Reproducible baselines} with 5-seed mean and standard
    deviation ($\sigma < 0.003$ F1) across 6 primary HMC benchmarks,
    establishing a new standard for statistical rigor in HMC evaluation.
\end{itemize}
